% ============================================================
% Integrante: Heily Rodríguez Rodríguez
% Fragmento puro: se incluye con \include desde SemArq1_TeamBales.tex.
% Fuente de trabajo: solucion.md (misma carpeta).
% Citas: páginas impresas de cada libro. El artículo web de Cockburn
% no tiene paginación y se cita por sección con \citep[«Sección»]{cockburn2005}.
% ============================================================

\chapter{P1: Trade-offs N-Capas vs. Clean/Onion/Hexagonal}

\section{El error fundamental de acoplamiento}

El error está en la \textbf{dirección de las dependencias}. En una arquitectura N-Capas tradicional, las dependencias del código fuente siguen el flujo de control: la presentación llama al negocio y el negocio llama a los datos. Como para llamar a un módulo hay que nombrarlo en el código, el negocio termina dependiendo de la base de datos \citep[pp.~44--45]{martin2017}. Esto invierte lo que debería ocurrir. La lógica de negocio es la parte más valiosa y estable del sistema, pero queda atada a la persistencia, que es la que más cambia. Martin advierte que un componente que se espera volátil no debe ser objeto de dependencia de otro difícil de cambiar \fuente{martin2017}{120}, y que un componente estable debe ser también abstracto \fuente{martin2017}{126}. Cockburn identifica el mismo defecto y lo llama \emph{enredo} entre la lógica de negocio y la interacción con el exterior \citep[«Nature of the Solution»]{cockburn2005}.

El problema se agrava porque con las dependencias también viajan los datos. Los formatos propios de la base de datos, como las filas devueltas por una consulta, atraviesan las capas, cuando Martin advierte que estas estructuras no deben cruzar las fronteras hacia dentro \citep[pp.~204, 207]{martin2017}. De este modo, el esquema de la base de datos se convierte en el modelo de todo el sistema. Por eso un cambio en una tabla obliga a modificar el negocio y, a través de él, la interfaz de usuario, que es justo lo que le sucede al equipo del enunciado.

\section{Solución estructural: Regla de Dependencia y Puertos y Adaptadores}

La \textbf{Regla de Dependencia} de Martin establece que las dependencias del código fuente solo pueden apuntar hacia dentro, hacia las políticas de más alto nivel. Nada de lo declarado en un círculo externo puede ser nombrado desde uno interno \fuente{martin2017}{203}. Para lograrlo se aplica la inversión de dependencias: el núcleo declara una interfaz, por ejemplo \texttt{IInventarioRepository}, y la capa de infraestructura la implementa. En ejecución, el negocio sigue invocando a la persistencia, pero en el código es la persistencia la que depende del negocio \citep[pp.~45, 206]{martin2017}. La interfaz de acceso a datos queda en el componente de negocio, mientras que su implementación queda en el de base de datos, que pasa a ser un \emph{plugin} de las reglas de negocio \citep[pp.~46, 168]{martin2017}.

Los \textbf{Puertos y Adaptadores} de Cockburn llegan a la misma estructura con otro vocabulario. Sustituyen la asimetría arriba/abajo de las capas por una asimetría dentro/fuera. Un \emph{puerto} es una conversación con un propósito concreto, expresada como una API. Un \emph{adaptador} traduce esa API a una tecnología determinada, y un mismo puerto puede tener varios adaptadores: uno para SQL, otro para ficheros o una base de datos simulada en memoria \citep[«Nature of the Solution»]{cockburn2005}. La base de datos es un actor secundario, es decir, uno al que la aplicación conduce, y se conecta a través de un puerto secundario \citep[«The Left-Right Asymmetry»]{cockburn2005}. En los dos enfoques ocurre lo mismo: la interfaz de persistencia pertenece al núcleo y la infraestructura depende de él.

\section{Impacto en la testabilidad y en el cambio de SQL Server a MongoDB}

Respecto a la testabilidad, con esta estructura las reglas de negocio pueden probarse sin interfaz de usuario, sin base de datos y sin ningún otro elemento externo \fuente{martin2017}{202}, que es también el objetivo declarado del patrón hexagonal \citep[«Intent»]{cockburn2005}. Basta con inyectar un doble de prueba, como un \emph{stub}, un \emph{mock} o un repositorio en memoria, en lugar del adaptador real. La guía de Microsoft describe esta práctica con repositorios falsos \fuente{delatorre2010}{146}, y Martin señala que los casos de uso son testables porque sus pasarelas a la base de datos pueden sustituirse por \emph{stubs} \fuente{martin2017}{214}. El resultado es que las pruebas dejan de ser lentas, porque evitan recursos costosos como la base de datos \fuente{martin2017}{252}, y dejan de ser frágiles, porque siguen la regla de no depender de cosas volátiles \fuente{martin2017}{251}.

En cuanto al cambio de SQL Server a MongoDB, basta con escribir un nuevo adaptador que implemente el mismo puerto y registrarlo en lugar del anterior. El dominio y los casos de uso no se modifican. Martin menciona este caso entre las propiedades comunes a estas arquitecturas, entre ellas la hexagonal: se puede cambiar SQL Server por Mongo u otra base de datos sin que las reglas de negocio queden atadas a ella \fuente{martin2017}{202}. La guía de Microsoft coincide: si las interfaces de los repositorios se definen en el dominio, la persistencia puede sustituirse sin afectarlo \citep[pp.~145--146]{delatorre2010}. Para que esto funcione, las entidades deben ser ignorantes de la persistencia y por la frontera solo deben cruzar estructuras simples \parencites[p.~207]{martin2017}[p.~55]{delatorre2010}. Conviene añadir un matiz honesto: Martin advierte que estos reemplazos pueden no ser triviales, aunque la estructura de \emph{plugins} los hace factibles \fuente{martin2017}{171}.



\section*{Notas de investigación}
% ============================================================
% Integrante: Heily Rodríguez Rodríguez
% Fragmento puro: se incluye desde preguntas/p1-trade-offs/solucion.tex
% mediante % ============================================================
% Integrante: Heily Rodríguez Rodríguez
% Fragmento puro: se incluye desde preguntas/p1-trade-offs/solucion.tex
% mediante % ============================================================
% Integrante: Heily Rodríguez Rodríguez
% Fragmento puro: se incluye desde preguntas/p1-trade-offs/solucion.tex
% mediante \input{investigacion/p1-regla-dependencia/notas}.
% Fuente de trabajo: notas.md (misma carpeta).
% Citas: páginas impresas de cada libro. El artículo web de Cockburn
% no tiene paginación y se cita por sección con \citep[«Sección»]{cockburn2005}.
% ============================================================

\section*{Notas de investigación: Regla de Dependencia y Puertos \& Adaptadores}

\section{Síntesis}

En la Arquitectura N-Capas tradicional, las dependencias del código fuente siguen el flujo de control. La presentación invoca al negocio y el negocio invoca a los datos, y para invocar a un módulo hay que nombrarlo \citep[pp.~44--45]{martin2017}. La propia guía de Microsoft describe esta relación tradicional, en la que las capas de alto nivel dependen de las de bajo nivel \fuente{delatorre2010}{64}. El resultado es que la lógica de negocio, que debería ser la parte más estable del sistema, conoce a la persistencia y depende de ella, aunque la persistencia es la parte más volátil. Martin pone como ejemplo de componente rígido precisamente el esquema de base de datos: concreto, volátil y del que dependen muchos otros componentes \fuente{martin2017}{129}. Cockburn identifica el mismo defecto y lo llama \emph{enredo} entre la lógica de negocio y la interacción con entidades externas. Además, observa que en la práctica no se respetan las líneas de los diagramas de capas \citep[«Nature of the Solution»]{cockburn2005}.

La \textbf{Regla de Dependencia} de Martin establece que las dependencias del código fuente solo pueden apuntar hacia dentro, hacia las políticas de más alto nivel: nada de un círculo interno puede nombrar algo declarado en uno externo \fuente{martin2017}{203}. Esto es posible gracias a la inversión de dependencias. Al interponer una interfaz, la dependencia del código queda orientada \emph{en contra} del flujo de control \citep[pp.~45, 206]{martin2017}. Así, la base de datos y la interfaz de usuario pasan a ser \emph{plugins} de las reglas de negocio \citep[pp.~46, 170\textendash 171]{martin2017}.

\textbf{Puertos y Adaptadores} llega a la misma estructura desde otro ángulo. La asimetría que importa no es arriba/abajo ni izquierda/derecha, sino dentro/fuera. Un \emph{puerto} es una conversación con un propósito, expresada como una API. Un \emph{adaptador} traduce esa API a una tecnología concreta, y un mismo puerto admite varios adaptadores: SQL, ficheros planos o una base de datos simulada en memoria \citep[«Nature of the Solution»]{cockburn2005}.

Para la P1, ambas propuestas dan el mismo beneficio. La lógica de negocio se puede probar sin interfaz de usuario ni base de datos \parencites[p.~202]{martin2017}[«Intent»]{cockburn2005}, y cambiar el motor de persistencia se reduce a escribir un nuevo adaptador. \citet[p.~202]{martin2017} menciona de forma explícita el ejemplo de sustituir SQL Server por Mongo. Hay un matiz importante: el propio Martin advierte que estos reemplazos pueden no ser triviales, aunque la estructura de \emph{plugins} los hace factibles \fuente{martin2017}{171}.

\section{Conceptos clave}

Los conceptos se agrupan en tres tablas: la Regla de Dependencia y la inversión de dependencias (Tabla~\ref{tab:inv-p1-regla}), la N-Capas tradicional y las métricas de acoplamiento (Tabla~\ref{tab:inv-p1-metricas}) y los Puertos y Adaptadores junto con la testabilidad (Tabla~\ref{tab:inv-p1-puertos}). Las definiciones son síntesis propias.

\begin{table}[htbp]
  \centering
  \small
  \caption{Conceptos clave de INV-P1: Regla de Dependencia e inversión de dependencias}
  \label{tab:inv-p1-regla}
  \begin{tabular}{@{}p{0.2\linewidth}p{0.51\linewidth}p{0.22\linewidth}@{}}
    \toprule
    Concepto & Definición & Fuente \\
    \midrule
    \raggedright Regla de Dependencia & Las dependencias del código fuente solo apuntan hacia dentro, hacia políticas de mayor nivel. Ningún nombre declarado en un círculo externo puede aparecer en uno interno. & \raggedright \fuente{martin2017}{203} \tabularnewline
    \raggedright Formatos de datos en la frontera & Dentro no deben usarse los formatos de datos de un círculo externo. Por la frontera solo cruzan estructuras simples, nunca filas de base de datos ni entidades. & \raggedright \citep[pp.~204, 207]{martin2017} \tabularnewline
    \raggedright Interface Adapters & Capa que convierte los datos entre el formato de los casos de uso y el de la base de datos o la web. Todo el SQL queda confinado en ella. & \raggedright \fuente{martin2017}{205} \tabularnewline
    \raggedright Dependencias frente a flujo de control & Sin polimorfismo, las dependencias del código siguen el flujo de control. Con una interfaz intermedia, cualquier dependencia puede invertirse. & \raggedright \citep[pp.~44--46]{martin2017} \tabularnewline
    \raggedright Cruce de fronteras & Cuando el control va de dentro hacia fuera, el círculo interno llama a una interfaz propia (\emph{output port}) que implementa el externo. & \raggedright \fuente{martin2017}{206} \tabularnewline
    \raggedright DIP & Los sistemas más flexibles dependen solo de abstracciones. Se evitan las concreciones \emph{volátiles}; las estables se toleran. & \raggedright \citep[pp.~87--88]{martin2017} \tabularnewline
    \raggedright Componente \texttt{main} & Las violaciones inevitables del DIP se concentran en pocos componentes concretos, típicamente \texttt{main}, donde se instancian las implementaciones. & \raggedright \fuente{martin2017}{91} \tabularnewline
    \raggedright Arquitectura de \emph{plugins} & Los componentes de base de datos e interfaz gráfica dependen de las reglas de negocio. Las interfaces de base de datos viven en el componente de negocio. & \raggedright \citep[pp.~168--171]{martin2017} \tabularnewline
    \bottomrule
  \end{tabular}
\end{table}

\begin{table}[htbp]
  \centering
  \small
  \caption{Conceptos clave de INV-P1: la N-Capas tradicional y las métricas de acoplamiento}
  \label{tab:inv-p1-metricas}
  \begin{tabular}{@{}p{0.2\linewidth}p{0.51\linewidth}p{0.22\linewidth}@{}}
    \toprule
    Concepto & Definición & Fuente \\
    \midrule
    \raggedright Capas cerradas & Cada petición atraviesa la capa inmediatamente inferior, lo que aísla los cambios, pero las dependencias siguen apuntando hacia la persistencia. & \raggedright \parencites[pp.~3--4]{richards2015}[p.~64]{delatorre2010} \tabularnewline
    \raggedright \emph{Fan-in} / \emph{Fan-out} & Dependencias entrantes y salientes de un componente. Corresponden al acoplamiento aferente ($C_a$) y eferente ($C_e$). & \raggedright \fuente{martin2017}{122} \tabularnewline
    \raggedright Inestabilidad ($I$) & $I = \textit{Fan-out} / (\textit{Fan-in} + \textit{Fan-out})$, con $I \in [0, 1]$. $I = 0$ es máximamente estable e $I = 1$, máximamente inestable. & \raggedright \parencites[p.~122]{martin2017}[cap.~3]{richardsford2020} \tabularnewline
    \raggedright Aferente y eferente & Términos de \emph{Structured Design} (Yourdon y Constantine, 1979): conexiones entrantes y salientes. & \raggedright \citep[cap.~3, «Coupling»]{richardsford2020} \tabularnewline
    \raggedright SDP & \emph{Depend in the direction of stability}: $I$ debe decrecer en la dirección de las dependencias. Una violación se corrige con el DIP. & \raggedright \citep[pp.~120, 123\textendash 125]{martin2017} \tabularnewline
    \raggedright SAP & Un componente debe ser tan abstracto como estable. La política de alto nivel va en componentes estables y abstractos. & \raggedright \fuente{martin2017}{126} \tabularnewline
    \raggedright Zona de dolor & Componentes muy estables y muy concretos: rígidos, difíciles de extender y de cambiar. El ejemplo de Martin es el esquema de base de datos. & \raggedright \fuente{martin2017}{129} \tabularnewline
    \bottomrule
  \end{tabular}
\end{table}

\begin{table}[htbp]
  \centering
  \small
  \caption{Conceptos clave de INV-P1: Puertos y Adaptadores y testabilidad}
  \label{tab:inv-p1-puertos}
  \begin{tabular}{@{}p{0.2\linewidth}p{0.51\linewidth}p{0.22\linewidth}@{}}
    \toprule
    Concepto & Definición & Fuente \\
    \midrule
    \raggedright Puerto & Conversación con un propósito entre la aplicación y el exterior, cuyo protocolo toma la forma de una API. & \raggedright \citep[«Nature of the Solution»]{cockburn2005} \tabularnewline
    \raggedright Adaptador & Traduce la API del puerto a una tecnología concreta. Puede haber varios por puerto. & \raggedright \citep[«Nature of the Solution»]{cockburn2005} \tabularnewline
    \raggedright Asimetría dentro/fuera & La asimetría que importa es interior/exterior. El hexágono elimina la lectura arriba-abajo de las capas apiladas. & \raggedright \citep[«Nature of the Solution»]{cockburn2005} \tabularnewline
    \raggedright Actores primarios y secundarios & El actor primario conduce la aplicación (\emph{driving}); el secundario es conducido por ella (\emph{driven}), como la base de datos. Puertos y adaptadores toman el nombre del actor al que sirven. & \raggedright \citep[«The Left-Right Asymmetry»]{cockburn2005} \tabularnewline
    \raggedright Base de datos simulada en memoria & Adaptador que sustituye a un actor secundario en las pruebas, como \texttt{RateRepository} / \texttt{MockRateRepository}, inyectado por constructor. & \raggedright \citep[«Sample Code»]{cockburn2005} \tabularnewline
    \raggedright Database Gateway & Interfaz polimórfica con las operaciones CRUD de la aplicación. Los casos de uso se prueban sustituyéndola por \emph{stubs}. & \raggedright \fuente{martin2017}{214} \tabularnewline
    \raggedright Repositorios falsos & Repositorios en memoria que sustituyen a los reales en las pruebas, configurados mediante el contenedor IoC. & \raggedright \citep[pp.~105, 146]{delatorre2010} \tabularnewline
    \raggedright DIP según la guía N-Capas & Los contratos residen en el alto nivel, así que el bajo nivel pasa a depender de él en tiempo de compilación. & \raggedright \citep[pp.~64--65]{delatorre2010} \tabularnewline
    \raggedright Contratos de repositorio en el dominio & Las interfaces de los repositorios se definen en el dominio y se implementan en infraestructura; las entidades ignoran la persistencia. & \raggedright \citep[pp.~55, 145\textendash 146]{delatorre2010} \tabularnewline
    \raggedright \emph{Fragile Tests Problem} & Las pruebas acopladas a partes volátiles se rompen en masa ante cambios triviales. Regla: \emph{Don't depend on volatile things}. & \raggedright \citep[pp.~251--252]{martin2017} \tabularnewline
    \bottomrule
  \end{tabular}
\end{table}

\section{Puntos de conexión con las preguntas}

\begin{itemize}
  \item \textbf{P1, error de acoplamiento.} Las dependencias del código siguen el flujo de control, así que el negocio depende de la persistencia \parencites[pp.~44--45]{martin2017}[p.~64]{delatorre2010}. Una capa volátil y concreta es objeto de dependencia de la lógica de negocio, en contra de la lectura cualitativa del SDP \fuente{martin2017}{120} y del SAP, lo que la sitúa en la zona de dolor \citep[pp.~126, 129]{martin2017}. Es el enredo entre negocio y exterior que describe \citet[«Nature of the Solution»]{cockburn2005}.
  \item \textbf{P1, solución estructural.} Se invierte la dependencia con interfaces definidas en el núcleo, siguiendo la Regla de Dependencia \citep[pp.~203, 206]{martin2017}. La base de datos pasa a ser un actor secundario, al que se accede por un puerto secundario mediante un adaptador intercambiable \citep[«Nature of the Solution», «The Left-Right Asymmetry»]{cockburn2005}.
  \item \textbf{P1, testabilidad.} Se sustituye el adaptador real por un \emph{mock} o un repositorio en memoria \parencites[«Sample Code»]{cockburn2005}[p.~214]{martin2017}[pp.~105, 146]{delatorre2010}. Las pruebas evitan recursos costosos como la base de datos \fuente{martin2017}{252} y dejan de depender de partes volátiles \fuente{martin2017}{251}.
  \item \textbf{P1, cambio a MongoDB.} Basta con un nuevo adaptador o una nueva implementación del repositorio \parencites[pp.~170--171, 202]{martin2017}[p.~146]{delatorre2010}, aunque el cambio no siempre es trivial \fuente{martin2017}{171}.
  \item \textbf{P2 y P4.} Para la P2: la asimetría dentro/fuera y la relación con el patrón \emph{Adapter} \citep[«Nature of the Solution», «Related Patterns»]{cockburn2005}. Para la P4: el DIP, el componente \texttt{main} y el \emph{output port} \citep[pp.~87--91, 206]{martin2017}.
\end{itemize}


.
% Fuente de trabajo: notas.md (misma carpeta).
% Citas: páginas impresas de cada libro. El artículo web de Cockburn
% no tiene paginación y se cita por sección con \citep[«Sección»]{cockburn2005}.
% ============================================================

\section*{Notas de investigación: Regla de Dependencia y Puertos \& Adaptadores}

\section{Síntesis}

En la Arquitectura N-Capas tradicional, las dependencias del código fuente siguen el flujo de control. La presentación invoca al negocio y el negocio invoca a los datos, y para invocar a un módulo hay que nombrarlo \citep[pp.~44--45]{martin2017}. La propia guía de Microsoft describe esta relación tradicional, en la que las capas de alto nivel dependen de las de bajo nivel \fuente{delatorre2010}{64}. El resultado es que la lógica de negocio, que debería ser la parte más estable del sistema, conoce a la persistencia y depende de ella, aunque la persistencia es la parte más volátil. Martin pone como ejemplo de componente rígido precisamente el esquema de base de datos: concreto, volátil y del que dependen muchos otros componentes \fuente{martin2017}{129}. Cockburn identifica el mismo defecto y lo llama \emph{enredo} entre la lógica de negocio y la interacción con entidades externas. Además, observa que en la práctica no se respetan las líneas de los diagramas de capas \citep[«Nature of the Solution»]{cockburn2005}.

La \textbf{Regla de Dependencia} de Martin establece que las dependencias del código fuente solo pueden apuntar hacia dentro, hacia las políticas de más alto nivel: nada de un círculo interno puede nombrar algo declarado en uno externo \fuente{martin2017}{203}. Esto es posible gracias a la inversión de dependencias. Al interponer una interfaz, la dependencia del código queda orientada \emph{en contra} del flujo de control \citep[pp.~45, 206]{martin2017}. Así, la base de datos y la interfaz de usuario pasan a ser \emph{plugins} de las reglas de negocio \citep[pp.~46, 170\textendash 171]{martin2017}.

\textbf{Puertos y Adaptadores} llega a la misma estructura desde otro ángulo. La asimetría que importa no es arriba/abajo ni izquierda/derecha, sino dentro/fuera. Un \emph{puerto} es una conversación con un propósito, expresada como una API. Un \emph{adaptador} traduce esa API a una tecnología concreta, y un mismo puerto admite varios adaptadores: SQL, ficheros planos o una base de datos simulada en memoria \citep[«Nature of the Solution»]{cockburn2005}.

Para la P1, ambas propuestas dan el mismo beneficio. La lógica de negocio se puede probar sin interfaz de usuario ni base de datos \parencites[p.~202]{martin2017}[«Intent»]{cockburn2005}, y cambiar el motor de persistencia se reduce a escribir un nuevo adaptador. \citet[p.~202]{martin2017} menciona de forma explícita el ejemplo de sustituir SQL Server por Mongo. Hay un matiz importante: el propio Martin advierte que estos reemplazos pueden no ser triviales, aunque la estructura de \emph{plugins} los hace factibles \fuente{martin2017}{171}.

\section{Conceptos clave}

Los conceptos se agrupan en tres tablas: la Regla de Dependencia y la inversión de dependencias (Tabla~\ref{tab:inv-p1-regla}), la N-Capas tradicional y las métricas de acoplamiento (Tabla~\ref{tab:inv-p1-metricas}) y los Puertos y Adaptadores junto con la testabilidad (Tabla~\ref{tab:inv-p1-puertos}). Las definiciones son síntesis propias.

\begin{table}[htbp]
  \centering
  \small
  \caption{Conceptos clave de INV-P1: Regla de Dependencia e inversión de dependencias}
  \label{tab:inv-p1-regla}
  \begin{tabular}{@{}p{0.2\linewidth}p{0.51\linewidth}p{0.22\linewidth}@{}}
    \toprule
    Concepto & Definición & Fuente \\
    \midrule
    \raggedright Regla de Dependencia & Las dependencias del código fuente solo apuntan hacia dentro, hacia políticas de mayor nivel. Ningún nombre declarado en un círculo externo puede aparecer en uno interno. & \raggedright \fuente{martin2017}{203} \tabularnewline
    \raggedright Formatos de datos en la frontera & Dentro no deben usarse los formatos de datos de un círculo externo. Por la frontera solo cruzan estructuras simples, nunca filas de base de datos ni entidades. & \raggedright \citep[pp.~204, 207]{martin2017} \tabularnewline
    \raggedright Interface Adapters & Capa que convierte los datos entre el formato de los casos de uso y el de la base de datos o la web. Todo el SQL queda confinado en ella. & \raggedright \fuente{martin2017}{205} \tabularnewline
    \raggedright Dependencias frente a flujo de control & Sin polimorfismo, las dependencias del código siguen el flujo de control. Con una interfaz intermedia, cualquier dependencia puede invertirse. & \raggedright \citep[pp.~44--46]{martin2017} \tabularnewline
    \raggedright Cruce de fronteras & Cuando el control va de dentro hacia fuera, el círculo interno llama a una interfaz propia (\emph{output port}) que implementa el externo. & \raggedright \fuente{martin2017}{206} \tabularnewline
    \raggedright DIP & Los sistemas más flexibles dependen solo de abstracciones. Se evitan las concreciones \emph{volátiles}; las estables se toleran. & \raggedright \citep[pp.~87--88]{martin2017} \tabularnewline
    \raggedright Componente \texttt{main} & Las violaciones inevitables del DIP se concentran en pocos componentes concretos, típicamente \texttt{main}, donde se instancian las implementaciones. & \raggedright \fuente{martin2017}{91} \tabularnewline
    \raggedright Arquitectura de \emph{plugins} & Los componentes de base de datos e interfaz gráfica dependen de las reglas de negocio. Las interfaces de base de datos viven en el componente de negocio. & \raggedright \citep[pp.~168--171]{martin2017} \tabularnewline
    \bottomrule
  \end{tabular}
\end{table}

\begin{table}[htbp]
  \centering
  \small
  \caption{Conceptos clave de INV-P1: la N-Capas tradicional y las métricas de acoplamiento}
  \label{tab:inv-p1-metricas}
  \begin{tabular}{@{}p{0.2\linewidth}p{0.51\linewidth}p{0.22\linewidth}@{}}
    \toprule
    Concepto & Definición & Fuente \\
    \midrule
    \raggedright Capas cerradas & Cada petición atraviesa la capa inmediatamente inferior, lo que aísla los cambios, pero las dependencias siguen apuntando hacia la persistencia. & \raggedright \parencites[pp.~3--4]{richards2015}[p.~64]{delatorre2010} \tabularnewline
    \raggedright \emph{Fan-in} / \emph{Fan-out} & Dependencias entrantes y salientes de un componente. Corresponden al acoplamiento aferente ($C_a$) y eferente ($C_e$). & \raggedright \fuente{martin2017}{122} \tabularnewline
    \raggedright Inestabilidad ($I$) & $I = \textit{Fan-out} / (\textit{Fan-in} + \textit{Fan-out})$, con $I \in [0, 1]$. $I = 0$ es máximamente estable e $I = 1$, máximamente inestable. & \raggedright \parencites[p.~122]{martin2017}[cap.~3]{richardsford2020} \tabularnewline
    \raggedright Aferente y eferente & Términos de \emph{Structured Design} (Yourdon y Constantine, 1979): conexiones entrantes y salientes. & \raggedright \citep[cap.~3, «Coupling»]{richardsford2020} \tabularnewline
    \raggedright SDP & \emph{Depend in the direction of stability}: $I$ debe decrecer en la dirección de las dependencias. Una violación se corrige con el DIP. & \raggedright \citep[pp.~120, 123\textendash 125]{martin2017} \tabularnewline
    \raggedright SAP & Un componente debe ser tan abstracto como estable. La política de alto nivel va en componentes estables y abstractos. & \raggedright \fuente{martin2017}{126} \tabularnewline
    \raggedright Zona de dolor & Componentes muy estables y muy concretos: rígidos, difíciles de extender y de cambiar. El ejemplo de Martin es el esquema de base de datos. & \raggedright \fuente{martin2017}{129} \tabularnewline
    \bottomrule
  \end{tabular}
\end{table}

\begin{table}[htbp]
  \centering
  \small
  \caption{Conceptos clave de INV-P1: Puertos y Adaptadores y testabilidad}
  \label{tab:inv-p1-puertos}
  \begin{tabular}{@{}p{0.2\linewidth}p{0.51\linewidth}p{0.22\linewidth}@{}}
    \toprule
    Concepto & Definición & Fuente \\
    \midrule
    \raggedright Puerto & Conversación con un propósito entre la aplicación y el exterior, cuyo protocolo toma la forma de una API. & \raggedright \citep[«Nature of the Solution»]{cockburn2005} \tabularnewline
    \raggedright Adaptador & Traduce la API del puerto a una tecnología concreta. Puede haber varios por puerto. & \raggedright \citep[«Nature of the Solution»]{cockburn2005} \tabularnewline
    \raggedright Asimetría dentro/fuera & La asimetría que importa es interior/exterior. El hexágono elimina la lectura arriba-abajo de las capas apiladas. & \raggedright \citep[«Nature of the Solution»]{cockburn2005} \tabularnewline
    \raggedright Actores primarios y secundarios & El actor primario conduce la aplicación (\emph{driving}); el secundario es conducido por ella (\emph{driven}), como la base de datos. Puertos y adaptadores toman el nombre del actor al que sirven. & \raggedright \citep[«The Left-Right Asymmetry»]{cockburn2005} \tabularnewline
    \raggedright Base de datos simulada en memoria & Adaptador que sustituye a un actor secundario en las pruebas, como \texttt{RateRepository} / \texttt{MockRateRepository}, inyectado por constructor. & \raggedright \citep[«Sample Code»]{cockburn2005} \tabularnewline
    \raggedright Database Gateway & Interfaz polimórfica con las operaciones CRUD de la aplicación. Los casos de uso se prueban sustituyéndola por \emph{stubs}. & \raggedright \fuente{martin2017}{214} \tabularnewline
    \raggedright Repositorios falsos & Repositorios en memoria que sustituyen a los reales en las pruebas, configurados mediante el contenedor IoC. & \raggedright \citep[pp.~105, 146]{delatorre2010} \tabularnewline
    \raggedright DIP según la guía N-Capas & Los contratos residen en el alto nivel, así que el bajo nivel pasa a depender de él en tiempo de compilación. & \raggedright \citep[pp.~64--65]{delatorre2010} \tabularnewline
    \raggedright Contratos de repositorio en el dominio & Las interfaces de los repositorios se definen en el dominio y se implementan en infraestructura; las entidades ignoran la persistencia. & \raggedright \citep[pp.~55, 145\textendash 146]{delatorre2010} \tabularnewline
    \raggedright \emph{Fragile Tests Problem} & Las pruebas acopladas a partes volátiles se rompen en masa ante cambios triviales. Regla: \emph{Don't depend on volatile things}. & \raggedright \citep[pp.~251--252]{martin2017} \tabularnewline
    \bottomrule
  \end{tabular}
\end{table}

\section{Puntos de conexión con las preguntas}

\begin{itemize}
  \item \textbf{P1, error de acoplamiento.} Las dependencias del código siguen el flujo de control, así que el negocio depende de la persistencia \parencites[pp.~44--45]{martin2017}[p.~64]{delatorre2010}. Una capa volátil y concreta es objeto de dependencia de la lógica de negocio, en contra de la lectura cualitativa del SDP \fuente{martin2017}{120} y del SAP, lo que la sitúa en la zona de dolor \citep[pp.~126, 129]{martin2017}. Es el enredo entre negocio y exterior que describe \citet[«Nature of the Solution»]{cockburn2005}.
  \item \textbf{P1, solución estructural.} Se invierte la dependencia con interfaces definidas en el núcleo, siguiendo la Regla de Dependencia \citep[pp.~203, 206]{martin2017}. La base de datos pasa a ser un actor secundario, al que se accede por un puerto secundario mediante un adaptador intercambiable \citep[«Nature of the Solution», «The Left-Right Asymmetry»]{cockburn2005}.
  \item \textbf{P1, testabilidad.} Se sustituye el adaptador real por un \emph{mock} o un repositorio en memoria \parencites[«Sample Code»]{cockburn2005}[p.~214]{martin2017}[pp.~105, 146]{delatorre2010}. Las pruebas evitan recursos costosos como la base de datos \fuente{martin2017}{252} y dejan de depender de partes volátiles \fuente{martin2017}{251}.
  \item \textbf{P1, cambio a MongoDB.} Basta con un nuevo adaptador o una nueva implementación del repositorio \parencites[pp.~170--171, 202]{martin2017}[p.~146]{delatorre2010}, aunque el cambio no siempre es trivial \fuente{martin2017}{171}.
  \item \textbf{P2 y P4.} Para la P2: la asimetría dentro/fuera y la relación con el patrón \emph{Adapter} \citep[«Nature of the Solution», «Related Patterns»]{cockburn2005}. Para la P4: el DIP, el componente \texttt{main} y el \emph{output port} \citep[pp.~87--91, 206]{martin2017}.
\end{itemize}


.
% Fuente de trabajo: notas.md (misma carpeta).
% Citas: páginas impresas de cada libro. El artículo web de Cockburn
% no tiene paginación y se cita por sección con \citep[«Sección»]{cockburn2005}.
% ============================================================

\section*{Notas de investigación: Regla de Dependencia y Puertos \& Adaptadores}

\section{Síntesis}

En la Arquitectura N-Capas tradicional, las dependencias del código fuente siguen el flujo de control. La presentación invoca al negocio y el negocio invoca a los datos, y para invocar a un módulo hay que nombrarlo \citep[pp.~44--45]{martin2017}. La propia guía de Microsoft describe esta relación tradicional, en la que las capas de alto nivel dependen de las de bajo nivel \fuente{delatorre2010}{64}. El resultado es que la lógica de negocio, que debería ser la parte más estable del sistema, conoce a la persistencia y depende de ella, aunque la persistencia es la parte más volátil. Martin pone como ejemplo de componente rígido precisamente el esquema de base de datos: concreto, volátil y del que dependen muchos otros componentes \fuente{martin2017}{129}. Cockburn identifica el mismo defecto y lo llama \emph{enredo} entre la lógica de negocio y la interacción con entidades externas. Además, observa que en la práctica no se respetan las líneas de los diagramas de capas \citep[«Nature of the Solution»]{cockburn2005}.

La \textbf{Regla de Dependencia} de Martin establece que las dependencias del código fuente solo pueden apuntar hacia dentro, hacia las políticas de más alto nivel: nada de un círculo interno puede nombrar algo declarado en uno externo \fuente{martin2017}{203}. Esto es posible gracias a la inversión de dependencias. Al interponer una interfaz, la dependencia del código queda orientada \emph{en contra} del flujo de control \citep[pp.~45, 206]{martin2017}. Así, la base de datos y la interfaz de usuario pasan a ser \emph{plugins} de las reglas de negocio \citep[pp.~46, 170\textendash 171]{martin2017}.

\textbf{Puertos y Adaptadores} llega a la misma estructura desde otro ángulo. La asimetría que importa no es arriba/abajo ni izquierda/derecha, sino dentro/fuera. Un \emph{puerto} es una conversación con un propósito, expresada como una API. Un \emph{adaptador} traduce esa API a una tecnología concreta, y un mismo puerto admite varios adaptadores: SQL, ficheros planos o una base de datos simulada en memoria \citep[«Nature of the Solution»]{cockburn2005}.

Para la P1, ambas propuestas dan el mismo beneficio. La lógica de negocio se puede probar sin interfaz de usuario ni base de datos \parencites[p.~202]{martin2017}[«Intent»]{cockburn2005}, y cambiar el motor de persistencia se reduce a escribir un nuevo adaptador. \citet[p.~202]{martin2017} menciona de forma explícita el ejemplo de sustituir SQL Server por Mongo. Hay un matiz importante: el propio Martin advierte que estos reemplazos pueden no ser triviales, aunque la estructura de \emph{plugins} los hace factibles \fuente{martin2017}{171}.

\section{Conceptos clave}

Los conceptos se agrupan en tres tablas: la Regla de Dependencia y la inversión de dependencias (Tabla~\ref{tab:inv-p1-regla}), la N-Capas tradicional y las métricas de acoplamiento (Tabla~\ref{tab:inv-p1-metricas}) y los Puertos y Adaptadores junto con la testabilidad (Tabla~\ref{tab:inv-p1-puertos}). Las definiciones son síntesis propias.

\begin{table}[htbp]
  \centering
  \small
  \caption{Conceptos clave de INV-P1: Regla de Dependencia e inversión de dependencias}
  \label{tab:inv-p1-regla}
  \begin{tabular}{@{}p{0.2\linewidth}p{0.51\linewidth}p{0.22\linewidth}@{}}
    \toprule
    Concepto & Definición & Fuente \\
    \midrule
    \raggedright Regla de Dependencia & Las dependencias del código fuente solo apuntan hacia dentro, hacia políticas de mayor nivel. Ningún nombre declarado en un círculo externo puede aparecer en uno interno. & \raggedright \fuente{martin2017}{203} \tabularnewline
    \raggedright Formatos de datos en la frontera & Dentro no deben usarse los formatos de datos de un círculo externo. Por la frontera solo cruzan estructuras simples, nunca filas de base de datos ni entidades. & \raggedright \citep[pp.~204, 207]{martin2017} \tabularnewline
    \raggedright Interface Adapters & Capa que convierte los datos entre el formato de los casos de uso y el de la base de datos o la web. Todo el SQL queda confinado en ella. & \raggedright \fuente{martin2017}{205} \tabularnewline
    \raggedright Dependencias frente a flujo de control & Sin polimorfismo, las dependencias del código siguen el flujo de control. Con una interfaz intermedia, cualquier dependencia puede invertirse. & \raggedright \citep[pp.~44--46]{martin2017} \tabularnewline
    \raggedright Cruce de fronteras & Cuando el control va de dentro hacia fuera, el círculo interno llama a una interfaz propia (\emph{output port}) que implementa el externo. & \raggedright \fuente{martin2017}{206} \tabularnewline
    \raggedright DIP & Los sistemas más flexibles dependen solo de abstracciones. Se evitan las concreciones \emph{volátiles}; las estables se toleran. & \raggedright \citep[pp.~87--88]{martin2017} \tabularnewline
    \raggedright Componente \texttt{main} & Las violaciones inevitables del DIP se concentran en pocos componentes concretos, típicamente \texttt{main}, donde se instancian las implementaciones. & \raggedright \fuente{martin2017}{91} \tabularnewline
    \raggedright Arquitectura de \emph{plugins} & Los componentes de base de datos e interfaz gráfica dependen de las reglas de negocio. Las interfaces de base de datos viven en el componente de negocio. & \raggedright \citep[pp.~168--171]{martin2017} \tabularnewline
    \bottomrule
  \end{tabular}
\end{table}

\begin{table}[htbp]
  \centering
  \small
  \caption{Conceptos clave de INV-P1: la N-Capas tradicional y las métricas de acoplamiento}
  \label{tab:inv-p1-metricas}
  \begin{tabular}{@{}p{0.2\linewidth}p{0.51\linewidth}p{0.22\linewidth}@{}}
    \toprule
    Concepto & Definición & Fuente \\
    \midrule
    \raggedright Capas cerradas & Cada petición atraviesa la capa inmediatamente inferior, lo que aísla los cambios, pero las dependencias siguen apuntando hacia la persistencia. & \raggedright \parencites[pp.~3--4]{richards2015}[p.~64]{delatorre2010} \tabularnewline
    \raggedright \emph{Fan-in} / \emph{Fan-out} & Dependencias entrantes y salientes de un componente. Corresponden al acoplamiento aferente ($C_a$) y eferente ($C_e$). & \raggedright \fuente{martin2017}{122} \tabularnewline
    \raggedright Inestabilidad ($I$) & $I = \textit{Fan-out} / (\textit{Fan-in} + \textit{Fan-out})$, con $I \in [0, 1]$. $I = 0$ es máximamente estable e $I = 1$, máximamente inestable. & \raggedright \parencites[p.~122]{martin2017}[cap.~3]{richardsford2020} \tabularnewline
    \raggedright Aferente y eferente & Términos de \emph{Structured Design} (Yourdon y Constantine, 1979): conexiones entrantes y salientes. & \raggedright \citep[cap.~3, «Coupling»]{richardsford2020} \tabularnewline
    \raggedright SDP & \emph{Depend in the direction of stability}: $I$ debe decrecer en la dirección de las dependencias. Una violación se corrige con el DIP. & \raggedright \citep[pp.~120, 123\textendash 125]{martin2017} \tabularnewline
    \raggedright SAP & Un componente debe ser tan abstracto como estable. La política de alto nivel va en componentes estables y abstractos. & \raggedright \fuente{martin2017}{126} \tabularnewline
    \raggedright Zona de dolor & Componentes muy estables y muy concretos: rígidos, difíciles de extender y de cambiar. El ejemplo de Martin es el esquema de base de datos. & \raggedright \fuente{martin2017}{129} \tabularnewline
    \bottomrule
  \end{tabular}
\end{table}

\begin{table}[htbp]
  \centering
  \small
  \caption{Conceptos clave de INV-P1: Puertos y Adaptadores y testabilidad}
  \label{tab:inv-p1-puertos}
  \begin{tabular}{@{}p{0.2\linewidth}p{0.51\linewidth}p{0.22\linewidth}@{}}
    \toprule
    Concepto & Definición & Fuente \\
    \midrule
    \raggedright Puerto & Conversación con un propósito entre la aplicación y el exterior, cuyo protocolo toma la forma de una API. & \raggedright \citep[«Nature of the Solution»]{cockburn2005} \tabularnewline
    \raggedright Adaptador & Traduce la API del puerto a una tecnología concreta. Puede haber varios por puerto. & \raggedright \citep[«Nature of the Solution»]{cockburn2005} \tabularnewline
    \raggedright Asimetría dentro/fuera & La asimetría que importa es interior/exterior. El hexágono elimina la lectura arriba-abajo de las capas apiladas. & \raggedright \citep[«Nature of the Solution»]{cockburn2005} \tabularnewline
    \raggedright Actores primarios y secundarios & El actor primario conduce la aplicación (\emph{driving}); el secundario es conducido por ella (\emph{driven}), como la base de datos. Puertos y adaptadores toman el nombre del actor al que sirven. & \raggedright \citep[«The Left-Right Asymmetry»]{cockburn2005} \tabularnewline
    \raggedright Base de datos simulada en memoria & Adaptador que sustituye a un actor secundario en las pruebas, como \texttt{RateRepository} / \texttt{MockRateRepository}, inyectado por constructor. & \raggedright \citep[«Sample Code»]{cockburn2005} \tabularnewline
    \raggedright Database Gateway & Interfaz polimórfica con las operaciones CRUD de la aplicación. Los casos de uso se prueban sustituyéndola por \emph{stubs}. & \raggedright \fuente{martin2017}{214} \tabularnewline
    \raggedright Repositorios falsos & Repositorios en memoria que sustituyen a los reales en las pruebas, configurados mediante el contenedor IoC. & \raggedright \citep[pp.~105, 146]{delatorre2010} \tabularnewline
    \raggedright DIP según la guía N-Capas & Los contratos residen en el alto nivel, así que el bajo nivel pasa a depender de él en tiempo de compilación. & \raggedright \citep[pp.~64--65]{delatorre2010} \tabularnewline
    \raggedright Contratos de repositorio en el dominio & Las interfaces de los repositorios se definen en el dominio y se implementan en infraestructura; las entidades ignoran la persistencia. & \raggedright \citep[pp.~55, 145\textendash 146]{delatorre2010} \tabularnewline
    \raggedright \emph{Fragile Tests Problem} & Las pruebas acopladas a partes volátiles se rompen en masa ante cambios triviales. Regla: \emph{Don't depend on volatile things}. & \raggedright \citep[pp.~251--252]{martin2017} \tabularnewline
    \bottomrule
  \end{tabular}
\end{table}

\section{Puntos de conexión con las preguntas}

\begin{itemize}
  \item \textbf{P1, error de acoplamiento.} Las dependencias del código siguen el flujo de control, así que el negocio depende de la persistencia \parencites[pp.~44--45]{martin2017}[p.~64]{delatorre2010}. Una capa volátil y concreta es objeto de dependencia de la lógica de negocio, en contra de la lectura cualitativa del SDP \fuente{martin2017}{120} y del SAP, lo que la sitúa en la zona de dolor \citep[pp.~126, 129]{martin2017}. Es el enredo entre negocio y exterior que describe \citet[«Nature of the Solution»]{cockburn2005}.
  \item \textbf{P1, solución estructural.} Se invierte la dependencia con interfaces definidas en el núcleo, siguiendo la Regla de Dependencia \citep[pp.~203, 206]{martin2017}. La base de datos pasa a ser un actor secundario, al que se accede por un puerto secundario mediante un adaptador intercambiable \citep[«Nature of the Solution», «The Left-Right Asymmetry»]{cockburn2005}.
  \item \textbf{P1, testabilidad.} Se sustituye el adaptador real por un \emph{mock} o un repositorio en memoria \parencites[«Sample Code»]{cockburn2005}[p.~214]{martin2017}[pp.~105, 146]{delatorre2010}. Las pruebas evitan recursos costosos como la base de datos \fuente{martin2017}{252} y dejan de depender de partes volátiles \fuente{martin2017}{251}.
  \item \textbf{P1, cambio a MongoDB.} Basta con un nuevo adaptador o una nueva implementación del repositorio \parencites[pp.~170--171, 202]{martin2017}[p.~146]{delatorre2010}, aunque el cambio no siempre es trivial \fuente{martin2017}{171}.
  \item \textbf{P2 y P4.} Para la P2: la asimetría dentro/fuera y la relación con el patrón \emph{Adapter} \citep[«Nature of the Solution», «Related Patterns»]{cockburn2005}. Para la P4: el DIP, el componente \texttt{main} y el \emph{output port} \citep[pp.~87--91, 206]{martin2017}.
\end{itemize}



